\section{Model Criticism and Nested Epistemic Self-Trust}

Gates 95--104 move from uncertainty \emph{within} a known model class to criticism of the class and the evidence-generating mechanism itself. This transition matters because ordinary posterior reweighting cannot repair a hypothesis space that omits the relevant explanation. The resulting architecture distinguishes belief about the world, belief about the model, evidence about model adequacy, and the control policy that decides when those layers license action.

\subsection{One-step policy self-correction}

Gate 95 gives an exact finite witness of a policy repairing its own overconfidence. The initial model belief is
\[
P(\mathrm{reliable})=\frac45,
\qquad
P(\mathrm{degraded})=\frac15.
\]
Under this belief the controller selects the fragile probe. The actual environment in the witness is degraded, where the objective probe values are
\[
V(\mathrm{fragile})=\frac1{10},
\qquad
V(\mathrm{safe})=\frac35.
\]
A failed calibration event has evidence $6/25$ and updates model belief to
\[
P(\mathrm{reliable}\mid\mathrm{fail})=\frac13,
\qquad
P(\mathrm{degraded}\mid\mathrm{fail})=\frac23.
\]
The revised policy selects the safe probe. Thus the verified witness contains a complete chain
\[
\text{overconfident model belief}
\longrightarrow
\text{calibration failure}
\longrightarrow
\text{model-posterior revision}
\longrightarrow
\text{action revision}.
\]

Gate 96 expresses the same correction through regret. Under the actual degraded model, switching from the fragile to the safe probe removes the finite decision loss encoded by the witness. This is a one-step self-correction result, not an almost-sure convergence theorem for repeated learning.

\subsection{When the hypothesis class itself is wrong}

Gate 97 adds an explicit outside candidate. The two familiar models predict positive rates
\[
\frac1{10}
\qquad\text{and}\qquad
\frac9{10},
\]
while the observed rate is $1/2$. Their squared discrepancies are both
\[
\left(\frac12-\frac1{10}\right)^2
=
\left(\frac12-\frac9{10}\right)^2
=
\frac4{25},
\]
which exceeds the encoded tolerance $1/10$. A supplied balanced model predicts $1/2$ exactly and therefore has discrepancy $0$.

The formal result is deliberately modest: expanding the model class \emph{can} repair an adequacy failure that no redistribution of mass inside the original class can fix. The missing model is supplied by the formalization; Gate 97 does not contain an automatic theory-formation algorithm.

Gate 98 makes the closed-class pathology explicit. A probability distribution can remain perfectly normalized while all models receiving that probability mass are misspecified. Normalization therefore cannot serve as a certificate of model adequacy.

\subsection{Surprise as a control signal for model criticism}

Gate 99 distinguishes routine evidence from an anomaly whose predictive evidence under the encoded closed-class model is
\[
P(\mathrm{anomaly})=\frac3{200}.
\]
This lies below the surprise threshold $1/20$ and triggers the meta-state \texttt{modelDoubt}. Gate 100 converts that meta-state into a structural action: routine evidence selects \texttt{keepClosed}, whereas the anomaly selects \texttt{expandUnknown}.

The distinction is operational. Ordinary Bayesian updating redistributes probability over the current hypotheses. Gate 100 instead changes which hypothesis space future inference may use. What it does \emph{not} do is synthesize the missing hypothesis. It formalizes the trigger for opening the model class.

\subsection{Unknown unknowns as an explicit epistemic gap}

Gate 101 refuses to force all probability back into the known class after model doubt. The closed open-model belief is
\[
(\ell,h,u)=\left(\frac12,\frac12,0\right),
\]
while the doubt state is
\[
(\ell,h,u)=\left(\frac14,\frac14,\frac12\right).
\]
Both are normalized. The known mass supports the proposition that the current model class is adequate; unknown mass supports its negation. Under the existing four-valued evidence semantics, the closed state is classified $T$, whereas the anomaly-induced open state is classified $N$.

This is not a claim that one half is a uniquely correct prior probability for unknown models. The amount is a finite design witness. The verified structural point is that probability can be reserved for outside-class possibility without pretending that a known model is adequate.

\subsection{Meta-level gaps and gluts}

Gate 102 separates lack of model knowledge from contradictory evidence about the model. Four explicit support pairs are classified as
\[
\left(\frac12,\frac12\right)\mapsto N,
\qquad
\left(\frac45,\frac45\right)\mapsto B,
\]
\[
\left(\frac9{10},\frac1{10}\right)\mapsto T,
\qquad
\left(\frac1{10},\frac9{10}\right)\mapsto F.
\]
Hence second-order ignorance and second-order contradiction remain extensionally distinct:
\[
N\neq B.
\]
An agent may therefore lack sufficient evidence about its own model, or possess strong evidence both for and against that model, without collapsing those epistemic defects into one generic uncertainty label.

\subsection{Two four-valued coordinates}

Gate 103 packages first-order and second-order assessment into
\[
\mathsf{NestedStatus}
=
\mathsf{FDEValue}_{\mathrm{world}}
\times
\mathsf{FDEValue}_{\mathrm{model}}.
\]
Explicit witnesses show the same world status $T$ alongside model status $T$, $N$, and $B$. Other witnesses hold the model coordinate fixed at $T$ while the world coordinate ranges over $T$, $F$, $B$, and $N$.

The formal claim is witness-level coordinate independence. Gate 103 does not yet establish that all sixteen pairs arise under one fixed dynamics. That stronger reachability question is deferred to Gate 105.

\subsection{Finite rational self-trust control}

Gate 104 turns the nested status into a deterministic controller with five actions:
\[
\mathsf{actPositive},\quad
\mathsf{actNegative},\quad
\mathsf{senseWorld},\quad
\mathsf{calibrateModel},\quad
\mathsf{reconcileModel}.
\]
Its priority is explicitly second-order. Model $N$ or $F$ triggers calibration; model $B$ triggers reconciliation. Only model $T$ licenses direct interpretation of the world coordinate. Under a trusted model, world $N$ or $B$ triggers additional world sensing, world $T$ licenses positive action, and world $F$ licenses negative action.

Thus an apparently decisive world verdict does not automatically license action when the mechanism responsible for that verdict is itself in doubt. The controller can represent the thought
\begin{quote}
The evidence currently points toward $p$, but I do not yet trust the process that produced that evidence enough to act on it.
\end{quote}
without replacing the world status by the model status.

\paragraph{Interpretation boundary.}
Gate 104 verifies the behavior of this particular finite priority policy. It does not prove that second-order maintenance must universally dominate first-order action. Gates 106 and 107 will make this limitation explicit by introducing epistemic costs and by showing that the qualitative four-valued compression can discard decision-relevant quantitative information.
